\section{Задача}

Сформировать варианты использования, разработать на их основе тестовое покрытие покрытие и провести функциональное тестирование интерфейса сайта (в соответствии с вариантом).

Вариант №5158: Разговоры с назиданием от моих друзей из Вселенной - http://angely-sveta.ru/

\begin{enumerate}
    \item Тестовое покрытие должно быть сформировано на основании набора прецедентов использования сайта.
    \item Тестирование должно осуществляться автоматически - с помощью системы автоматизированного тестирования Selenium.
    \item Шаблоны тестов должны формироваться при помощи Selenium IDE и исполняться при помощи Selenium RC в браузерах Firefox и Chrome.
    \item Предполагается, что тестируемый сайт использует динамическую генерацию элементов на странице, т.е. выбор элемента в DOM должен осуществляться не на основании его ID, а с помощью XPath.
\end{enumerate}

\section{Список сценариев использования}

\subsection{Общий путь пользователя}
\begin{enumerate}
    \item Доступ к домашней странице и навигация по основным разделам
    \item Переключение между языковыми версиями (если доступны)
    \item Чтение статей и духовного контента
    \item Просмотр изображений/медиа (если доступны)
    \item Доступ к контактной информации
\end{enumerate}

\subsection{Конкретные сценарии использования}
\begin{enumerate}
    \item Навигация по домашней странице — Пользователь посещает домашнюю страницу и получает доступ к главному меню
    \item Просмотр контента — Пользователь может читать статьи и духовный контент
    \item Просмотр медиа — Пользователь может просматривать изображения, видео или другие медиа
    \item Переключение языков — Пользователь может переключаться между языковыми версиями
    \item Контактная информация — Пользователь может найти контактные данные или форму обратной связи
\end{enumerate}

\section{Чеклист тестирования}

\begin{enumerate}
    \item Тестирование навигации по домашней странице
    \item Тестирование переключения языков
    \item Тестирование просмотра контента
    \item Тестирование просмотра медиа
    \item Тестирование контактной информации
\end{enumerate}

\section{Результаты тестирования}

Для автоматизированного тестирования был использован Playwright.
Тесты рапсположены в директории tests.
Для обеспечения корректного отображения и сохранения верстки был использован visual testing.
